\section{Measurement implications and discussion}
\paragraph{Observe required behavior rather than implementation bookkeeping.}
X05 and X28 expose a measurement problem: an implementation can return a contract relevant result without populating an optional ledger. Evaluators should observe required outputs and externally visible effects, using internal traces as supporting evidence. X06 illustrates a second distinction. Its mixed record witness establishes confidentiality without requiring useful public output. Appendix~\ref{app:observers} specifies prospective output rules separately from the historical results.

\begin{table}[h]
\centering\footnotesize
\begin{tabularx}{\linewidth}{@{}>{\raggedright\arraybackslash}p{2.9cm}>{\raggedright\arraybackslash}X>{\raggedright\arraybackslash}X@{}}
\toprule
Claim & Evidence & Alternative explanation retained\\
\midrule
Observed estimates differ in direction & C$-$N is positive, zero, and negative & Sampling variation or family cancellation can produce this pattern.\\
Reminder response varies by configuration & Codex contrasts and bounds agree in direction; MiniSWE does not & Interfaces and instructions may shape the response.\\
Visible statement differs from protection & Fifteen runs have protection without agreed prior assumption recognition & The codes do not identify memory use or a causal pathway.\\
Focal passes depend on the observer & X05/X28 corrections and separate X06 output requirements & Confidentiality can hold because delivery was rejected.\\
\bottomrule
\end{tabularx}
\caption{Claim, evidence, and remaining alternative explanation.}
\label{tab:claims}
\end{table}

\paragraph{Limitations.}
The synthetic tasks provide explicit source assumptions and separately testable functionality and focal properties; thirteen selected families with two repetitions offer limited precision and do not establish repository scale external validity. Their trust shifts cover a constrained part of procedural applicability. The exploratory behavioral analysis uses two model assisted passes, with eight disagreements and no human validation. Retrieval, naturally occurring histories, and longer term memory dynamics remain outside the supplied memory intervention.

\section{Conclusion}
Supplying correct procedural memory produced both unchecked reuse and successful adaptation in the saved coding trajectories. The aggregate endpoint also fell when agents completed fewer tasks, showing why a reduction in functional focal violations is insufficient on its own. A generic applicability reminder had a consistent negative direction in the three Codex configurations, including under the bounds for missing recorded outcomes; the MiniSWE configurations did not share that pattern. Evaluating procedural transfer requires observing task completion, the focal property, and the implementation that connects the two.
